% Historical proof of concept: single-run numbered citations.
%
% This file predates monoref v1.4.  Since v1.4 the package itself
% provides single-run \cite (multi-key lists, the optional note, the
% [label] \bibitem form, hyperref links, and a natbib/cite/biblatex
% guard) -- see examples/07-citations.tex in the distribution.
%
% The original three-line wrapper that proved the idea, kept for the
% record (v1.4 supersedes it; loading it alongside v1.4 is harmless,
% the package's definitions win at begin-document):
%
%   \let\origbibitem\bibitem
%   \def\bibitem#1{\origbibitem{#1}\label{bib:#1}}
%   \renewcommand\cite[1]{[\monoref@refinner{bib:#1}]}
%
% Compile ONCE:  lualatex bib-poc.tex
\documentclass{article}
\usepackage{monoref}

\begin{document}

\section{Text}\label{sec:text}
Forward citations before the bibliography exists: first~\cite{knuth},
then~\cite{lamport}, then~\cite{luatex}.  A repeat of the middle
one~\cite{lamport}.  Page of bibliography: \pageref{sec:bib}.

\begin{thebibliography}{9}\label{sec:bib}
\bibitem{knuth} D. Knuth, \emph{The \TeX book}, 1986.
\bibitem{lamport} L. Lamport, \emph{\LaTeX: A Document Preparation System}, 1994.
\bibitem{luatex} LuaTeX development team, \emph{LuaTeX reference manual}, 2024.
\end{thebibliography}

\end{document}
